@@M0@@. Iterating the operation @@M1@@ times, we reach @@M2@@, which is zero, and have therefore written @@M3@@ as an element of @@M4@@. In view of (6.6.5), the two filtrations are therefore identical.

\sgaStableTarget{sga6.e05.para.6-7}6.7. We now suppose that @@M5@@ is a @@M6@@-ring augmented over @@M7@@, the augmentation homomorphism being denoted by @@M8@@. We therefore have on @@M9@@ the canonical filtration of an augmented @@M10@@-ring, and can thus associate to @@M11@@ its Chern ring.

If @@M12@@, then @@M13@@ by definition of the filtration of an augmented @@M14@@-ring. We denote by @@M15@@, and call the \uline{total Chern class of} @@M16@@, the image of @@M17@@ in @@M18@@; we denote by @@M19@@, and call the @@M20@@-th Chern class of @@M21@@, the @@M22@@-th coefficient of @@M23@@; we therefore have:
\sgaStableTarget{sga6.e05.eq.6-7-1}@@M24@@
We then put:
\sgaStableTarget{sga6.e05.eq.6-7-2}@@M25@@
and call @@M26@@ the completed Chern class of @@M27@@.

\sgaStableTarget{sga6.e05.prop.6-8}\textbf{Proposition 6.8.} \uline{The completed Chern class} @@M28@@, \uline{is} \uline{a} @@M29@@-@@M30@@-\uline{homomorphism of} \uline{augmented} @@M31@@-@@M32@@-\uline{algebras}.

It is clear that the homomorphism thus defined is additive; to verify that it is multiplicative, we write:
@@M33@@
